\chapter{34}



Louis’ Frage schien Joh im ersten Moment zu erstaunen. Mit Blick in die
Ferne schüttelte er nun den Kopf. Dann lachte er.

«Ach,» er vergrösserte seine Augen, «willst du es wirklich wissen?»

«Ja. Schon gestern hab ich mich über euch gewundert. Ihr seid eine
gelungene Kombination,» Louis grinste und wunderte sich, wie ihn Johs
Offenheit anzustecken vermochte, «es würde mich echt interessieren, wie
das mit euch angefangen hat.»

«Hm,» Joh kratzte sich am Kinn und dann wirkte er, als glaube er selbst
nicht, was er erzählte, «mir ist, als wär’ es gestern passiert. Ich wurde
in die Rechtsmedizin in Zürich bestellt, um eine Leiche zu
identifizieren. Zu dieser Zeit arbeitete ich in der Stadt auf dem
Sozialamt. Es war ein schwerer Gang, einer unserer Jungs war bei einer
Messerstecherei ums Leben gekommen. Seine Eltern lebten in Südamerika,
waren nicht erreichbar, wie immer, und weil der Junge bei uns im
Programm war, lag es an mir, ihn zu identifizieren. Schon vom Geruch in
der Leichenhalle wurde mir schlecht. Und als ich Iker auf diesem kalten,
silbernen Tisch liegen sah, knickten mir fast die Beine weg. Ich riss
mich zusammen und höre mich noch heute «Ja, das ist Iker» sagen. Erst
jetzt nahm ich die Frau wirklich wahr, die nun das Leichentuch über
seinen Kopf hochzog. Sie fasste mich sanft am Arm und führte mich ein
paar Schritte vom Tisch weg. Und dann, Ciro,» Joh rang nach Worten und
fixierte Louis, «schaute sie mich an, wie noch nie jemand zuvor, schwer
in Worten auszudrücken, liebevoll, wissend, anteilnehmend, ich weiss
nicht.»

Louis stand bewegungslos da. Joh richtete sich auf und atmete schwer
aus. Dann schlug er sich die Hände an die Oberschenkel.

«Im Unglück des einen liegt manchmal das Glück des andern», dachte ich
mir später. Ja, so lernten wir uns kennen, Manuela und ich. Ihr ist
unsere erste Begegnung in ähnlicher Erinnerung geblieben. Wir wurden zu
einem Paar. Liessen es langsam angehen, vorsichtig auch. Beide waren wir
enttäuschte Liebende, könnte man sagen, aber wer schon nicht,» Joh hob
die Hände während er fortfuhr, «Ikers Schicksal gab mir den letzten
Kick, endlich meine Vision von einer anderen Form von Resozialisierung
zu verwirklichen; er ist eigentlich der stille Pate der
\emph{Sunnewaid}. Seinen eingeschnitzten Vornamen findest du draussen
über dem Eingang zur Küche.»

«Eine traurig-schöne Geschichte, wow,» sagte Louis leise und fühlte sich
Joh gerade nah. Da war aber auch gleichzeitig dieses undefinierbare
Gefühl. Warum gab Joh so viel Intimes preis? Wo war Louis hier gelandet?
Ihre beruflichen Hintergründe könnten ihm gefährlich werden. Er musste
auf der Hut sein.

«Erst später, als wir uns für die Übernahme und Bewirtschaftung des
Hofes entschlossen hatten, bildete sich Manuela in Tiermedizin weiter
und unsere Hofpraxis konnte über die Zeit entstehen, nachdem sie die
Zulassung bekommen hatte.»

Joh verschloss die Tür und führte Louis zurück. Beim Gang durch den
Vorraum zur Wohnung, drehte sich Louis um.

«Beeindruckend, was ihr euch hier aufgebaut habt, interessant.»

Johs Blick verriet, wie gut er sich gerade fühlte.

Zurück in der Küche spürte Louis den Drang, sich endlich im Internet
nach Giorgias Schicksal umzusehen. Er brauchte Zeit. Alleine für sich.
Johs Vorschlag kam gelegen.

«Nun denn, Ciro,» Joh hielt Louis freundschaftlich an der Schulter, «das
wär’s mal vorläufig.»

Jetzt sah Joh sehr ernst aus.

«Vielleicht ein bisschen viel für den Anfang, Ciro. Doch unser oberstes
Credo ist Offenheit und Transparenz. Das erwarten wir auch von unseren
Mitarbeitenden. Nun denn, ich räum die Küche auf und dann können wir uns
über eine mögliche Anstellung unterhalten, okay, in einer Stunde?»

Louis nickte. An Offenheit fehlte es den beiden wahrlich nicht.

«Sehr gern, ja.»

\sectionbreak

Erleichtert stieg Louis die Treppe hoch. Dieses kleine Timeout brauchte
er nun dringend.

Sein Herz klopfte heftiger, als er sich auf sein Bett setzte und das
Handy entriegelte.

Erstmal Recherche zu Giorgias Zustand. Alles andere später.

Er schwankte, wusste nicht, sollte er sich auf irgendwelche
Informationen über Giorgias Schicksal freuen oder wäre es besser, nichts
zu finden.
